%%==============================================================================
%% 中文摘要 / 英文摘要
%%==============================================================================

\begin{cabstract}
大语言模型在多类应用中得到广泛使用，其生成内容的安全性与模型能力保持已成为重要问题。参数级知识编辑通过对模型局部参数进行干预，能够在不进行全量重训练的情况下改变特定知识或行为，为模型发布后的安全维护提供了轻量途径。然而，编辑结果可能受到后续微调影响；同时，不同风险类别和等级对模型行为的要求并不相同，统一编辑容易造成安全防护不足或正常请求过度拒答。针对这些问题，本文围绕安全知识编辑的微调鲁棒性、多类别风险控制及其系统应用开展研究。

（1）针对编辑结果在后续微调中可能衰减的问题，研究面向微调鲁棒性的双子空间约束安全知识编辑方法。分别依据安全编辑目标与正常任务样本构建安全子空间和效用保护子空间，并将二者共同纳入参数更新约束；进一步对后续微调造成的参数扰动进行建模，通过鲁棒优化提升编辑在不同微调条件下的保持能力。研究将从安全编辑效果、微调后保持率、通用能力变化及编辑局部性等方面评价方法。

（2）针对多类别有害内容编辑中的类别干扰与风险处置粒度不足问题，研究面向多类别安全性的分级解耦与动态路由编辑方法。该方法按照风险类别和风险等级组织相对独立的编辑模块，通过解耦约束减轻联合编辑时的跨类别干扰，并依据输入内容的风险特征动态选择或组合相应模块，以实现差异化安全干预。研究将重点考察类别编辑效果、模块隔离能力、路由性能、正常请求效用及过度拒答情况。

（3）面向方法的工程应用，设计大语言模型生成内容安全增强系统，集成基础模型管理、鲁棒安全编辑、多类别安全编辑、风险感知动态路由及安全评估等功能，形成从编辑配置到效果评估的系统流程。

本文从编辑持久性和编辑粒度两个角度研究参数级安全知识编辑，并通过实验设计与系统实现考察方法的有效性、适用边界及工程可用性，为大语言模型内容安全的精细化维护提供参考。

\ckeywords{大语言模型；生成内容安全；知识编辑；微调鲁棒性；动态路由}
\end{cabstract}

\begin{eabstract}
Large language models are widely used in a range of applications, making the safety of generated content and the preservation of model capabilities important concerns. Parameter-level knowledge editing intervenes in localized model parameters to change selected knowledge or behavior without full retraining, offering a lightweight approach to post-deployment safety maintenance. However, subsequent fine-tuning may weaken an edit. Moreover, different risk categories and levels call for different model behaviors, while uniform editing can either leave safety gaps or cause excessive refusals of benign requests. This thesis studies fine-tuning robustness, multi-category risk control, and system support for safety-oriented knowledge editing.

(1) To address the potential decay of edits after subsequent fine-tuning, this thesis studies a dual-subspace-constrained safety editing method for fine-tuning robustness. A safety subspace is constructed from safety-editing objectives, and a utility-preservation subspace is constructed from benign task examples. Both subspaces are incorporated into parameter-update constraints. Perturbations caused by subsequent fine-tuning are further modeled, and robust optimization is used to improve edit retention under different fine-tuning conditions. The method is evaluated in terms of safety-editing effectiveness, post-fine-tuning retention, general capability changes, and edit locality.

(2) To address cross-category interference and insufficient risk granularity in editing multiple types of harmful content, this thesis studies a graded, decoupled, and dynamically routed editing method for multi-category safety. Relatively independent editing modules are organized by risk category and level. Decoupling constraints are used to reduce interference during joint editing, while input risk features guide the dynamic selection or combination of modules for differentiated safety interventions. The evaluation focuses on category-level editing effectiveness, module isolation, routing performance, utility on benign requests, and over-refusal.

(3) To support practical use, a safety-enhancement system for LLM-generated content is designed. It integrates base-model management, robust safety editing, multi-category editing, risk-aware dynamic routing, and safety evaluation into a workflow from edit configuration to outcome assessment.

By addressing edit persistence and editing granularity, this thesis investigates parameter-level safety knowledge editing and examines its effectiveness, applicability, and engineering feasibility through an experimental design and system implementation. The work provides a reference for fine-grained maintenance of LLM content safety.

\ekeywords{Large Language Models; Generated Content Safety; Knowledge Editing; Fine-Tuning Robustness; Dynamic Routing}
\end{eabstract}
